\documentclass[a4paper,10pt]{article}

\usepackage{fancyhdr}
\usepackage{graphicx}
\usepackage{geometry}
\geometry{a4paper, left=2cm, right=2cm, top=1.5cm, bottom=3cm }
\usepackage{caption}
\usepackage{subcaption}
\usepackage{hyperref}
\usepackage{natbib}

\usepackage{etoolbox,fancyhdr,xcolor}
\newcommand{\headrulecolor}[1]{\patchcmd{\headrule}{\hrule}{\color{#1}\hrule}{}{}}
\newcommand{\footrulecolor}[1]{\patchcmd{\footrule}{\hrule}{\color{#1}\hrule}{}{}}
\renewcommand{\headrulewidth}{1pt}
\headrulecolor{red!100}%
\renewcommand{\footrulewidth}{1pt}
\footrulecolor{red!100}%

\fancyhf{}
\fancyhead[R]{\includegraphics[width=0.25\textwidth]{NCSUImage.png}}

\fancyfoot[L]{ST 704 Final Project}
\fancyfoot[C]{Department of Statistics}
\fancyfoot[R]{\thepage}

\setlength{\headheight}{15mm}
\pagestyle{fancy}

\bibliographystyle{apalike}

\usepackage{times}
\begin{document}

\noindent 
\begin{center}
\textbf{{\Large ST 704 Final Project Interim Report}} \\
\end{center}

\noindent 
\textbf{Carter Hall, Connor McNeill, Miles Woollacott} \textit{North Carolina State University}
\\

\noindent
This report is a quick overview of the work that has been completed, as well as where the group plans to apply these efforts.
\begin{enumerate}
    \item \textbf{Data Cleaning}: With multiple data sources within the repository we identified as our dataset for this project, a natural step before conducting any analysis and/or modeling is to clean and preprocess the data. 

    This task has proved arduous with thousands of school names across the United States, hundreds of which have alternative names, have \textit{changed} names, or have different naming conventions between the two datasets. This has culminated in a laborious effort to identify and \textit{hand-change} school names so that joining the datasets together introduces no unnecessary rows (i.e., NA values).  
    \item \textbf{Exploratory Data Analysis}: Given that we could identify \textit{multiple} response variables to model in explaining our research question, our EDA has been targeted at familiarizing ourselves with the data and identifying surprising [linear] associations. 
    \item \textbf{Identification of New Modeling Techniques}: Rather than simply \textit{re}-applying regression techniques from a Bayesian paradigm, we floated the possibility of tree-based methods (e.g., \textbf{Bayesian Additive Regression Trees} (BART), classification/regression trees, etc.) as a new modeling technique. Once the data cleaning is \textit{complete} (which it would be at the deadline of this report), group discussions will certainly pin down which of these techniques is finalized. Separately, \textit{regression splines} were an idea we hypothesized, including methods like \textbf{Multivariate additive regression spline} (MARS). \textit{The working plan for modeling is to fit a plethora of models and identify those which are performative, parsimonous, and as interpretable as possible}! 
\end{enumerate}
\end{document}